% Checkpoint37. Direct GT0 construction; no assumed three-strand GT0 action.
\statement{A.130}{Proved}{The actual GT0 two-strand reflection and its noncusp fixed group}
Let \(\alpha=(-1,\phi)\) be the original normalized negative
involution in \(\widehat{\mathrm{GT}}_0\), and let
\[
 E=\widehat F(X,Y),\quad K=[Y,X],\quad
 z=(KX)^{-1}=YX^{-1}Y^{-1},\quad
 \Phi=\phi(X,K),\quad Q=\phi(z,K).
\]
The generator assignments
\[
 T(X)=\Phi X^{-1}\Phi^{-1},\qquad
 T(Y)=K^{-1}QY\Phi^{-1}                                     \tag{A.130.1}
\]
define an actual involutory automorphism of \(E\), with
\[
 T(K)=K^{-1},\qquad T(\Phi)=\Phi^{-1},\qquad Tj=j^{-1}T,       \tag{A.130.2}
\]
where \(j\) is the original two-strand action of A.126.
Put
\[
 \beta=\operatorname{Ad}(\Phi^{-1})T,\qquad
 \theta=\operatorname{Ad}(Y^{-1})j,\qquad
 b=\operatorname{Ad}(X)\beta,\qquad
 t=\operatorname{Ad}(X)\theta .
                                                               \tag{A.130.3}
\]
Then \(\beta,\theta,b,t\) are involutions, \(\beta\theta=\theta\beta\)
and \(bt=tb\). The actual noncusp fixed group is
\[
          E^b=\overline{\langle a_B\rangle}
                \simeq\widehat{\mathbb Z},\qquad
          (a_B)_{\rm ab}=[Y],                               \tag{A.130.4}
\]
with uniquely determined \(a_B\). Every noncusp section centralizer
for this actual outer action is conjugate to this group.
An exact \(\theta\)-centralizing lift is unique within its outer
class. No full three-strand \(\widehat{\mathrm{GT}}_0\) elliptic
action, canonical-word pentagon at that scope, or full generation
conclusion is asserted.

\begin{proof}
\emph{1. Both ordered original hexagons.}
The defining negative two-cycle and three-cycle, in the exact
normalization of C.30, are
\[
 \phi(x,y)\phi(y,x)=1,\qquad
 \phi(x,y)x^{-1}\phi(z_0,x)z_0^{-1}\phi(y,z_0)y^{-1}=1,
 \quad z_0=(xy)^{-1}.                                      \tag{A.130.5}
\]
The same word satisfies the \(m=0\) three-cycle as well.
To derive it directly, put
\(A_0=\phi(x,y)\), \(B_0=\phi(y,z_0)\),
\(C_0=\phi(z_0,x)\), and apply the second identity of
(A.130.5) under \((x,y)\mapsto(x,z_0)\).
The new third argument is \(z_0^{-1}yz_0=xyx^{-1}\).
The two-cycle and conjugation equivariance give the same
substitution as in the complete C.30 proof, namely
\(C_0^{-1}A_0^{-1}B_0^{-1}=1\).
Thus \(A_0C_0B_0=1\), without assuming a further braid action.
At \(x=K,y=X\) this yields
\[
 \phi(X,z)=Q^{-1}\Phi,\qquad
                Qz^{-1}Q^{-1}=K\Phi X\Phi^{-1},              \tag{A.130.6}
\]
where the second equality is the negative hexagon in
(A.130.5) with the first equality substituted.

\emph{2. The actual automorphism and involution.}
The assignments (A.130.1) first define a continuous endomorphism.
Its commutator image is, in the original order,
\[
 \begin{aligned}
 T(K)&=K^{-1}QYX^{-1}Y^{-1}Q^{-1}K\Phi X\Phi^{-1}\\
     &=K^{-1}QzQ^{-1}K\Phi X\Phi^{-1}=K^{-1}.
 \end{aligned}                                              \tag{A.130.7}
\]
The last equality is the inverse of (A.130.6).
The general finite-cover lemma proved in A.123 now makes \(T\)
an automorphism. Explicitly, a proper closed image lies in a
proper open subgroup containing \(K\). In its finite
punctured-torus cover the distinguished boundary has degree one,
and there is at least one other boundary. Its homology class is
primitive and nonzero, whereas (A.130.7) would make it a
commutator of elements of that open subgroup. This contradiction
proves surjectivity; finite generation and Hopficity give
automorphy. No maximal pro-\(p\) generation argument is substituted.

The subgroup \(H=\overline{\langle K,X\rangle}\) is the actual
free profinite factor described in A.122: the open subgroup
\(\ker(\exp_Y\bmod2)\) has basis
\(X,YXY^{-1},Y^2\), and replacing \(YXY^{-1}\) by
\(K=(YXY^{-1})X^{-1}\) gives the assertion. The restriction
of \(T\) to \(H\) is exactly the original \(\alpha\) under
\((x,y)=(K,X)\), since \(\phi(K,X)=\Phi^{-1}\).
The original involution identity therefore gives
\(T(\Phi)=\Phi^{-1}\).
It follows that \(T^2\) fixes \(X,K\); its abelianization is
the identity because \(T_{\rm ab}=\operatorname{diag}(-1,1)\).
An automorphism fixing \(X,K\) has \(Y\)-image \(YX^r\):
the equality of the two commutators implies that
\(Y^{-1}f(Y)\) centralizes \(X\), and C.5 gives
\(C_E(X)=\overline{\langle X\rangle}\).
Identity abelianization forces \(r=0\). Hence \(T^2=1\).

\emph{3. The full half-twist and commuting normalization.}
We have \(Q=j(\Phi)\), \(j^{-1}\Phi=K^{-1}QK\) and
\(j^{-1}(X^{-1})=XK\). Both \(Tj\) and \(j^{-1}T\)
send \(K\) to \(K^{-1}\) and send \(X\) to
\[
          K^{-1}Qz^{-1}Q^{-1}K=\Phi X\Phi^{-1}K .
\]
Their abelianizations agree. The same \(X,K\) argument proves
\(Tj=j^{-1}T\). The equality \(T(\Phi)=\Phi^{-1}\) gives
\(\beta^2=1\), and A.126 gives \(\theta^2=1\).
To check their commutation directly, use
\[
 \beta(X)=X^{-1},\quad
 \beta(Y)=\Phi^{-1}K^{-1}Y\theta(\Phi),\quad
 \beta(K)=\Phi^{-1}K^{-1}\Phi,\quad
 \theta(K)=Y^{-1}KY .
\]
Both \(\beta\theta\) and \(\theta\beta\) send \(X\) to \(X\)
and send \(K\) to
\(\theta(\Phi)^{-1}Y^{-1}K^{-1}Y\theta(\Phi)\).
Their abelianizations agree, so they are equal by the same
argument. Since \(\beta(X)=\theta(X)=X^{-1}\), (A.130.3)
now also gives \(b^2=t^2=1\) and \(bt=tb\).

\emph{4. Actual finite-factor centralizers and rigidity.}
In the free basis \(X,Z=YX^{-1}\), \(\theta\) inverts both
generators. Thus
\[
 E\rtimes\langle\theta\rangle
     =C_2(\theta)\amalg C_2(X\theta)\amalg C_2(Z\theta).
                                                               \tag{A.130.8}
\]
The inverse maps use \(X=(X\theta)\theta\) and
\(Z=(Z\theta)\theta\).
The finite-factor centralizer theorem already applied in A.119
\citep[Theorem~B']{A_HR1985} therefore gives
\(E^\theta=E^t=1\).
If \(\operatorname{Ad}(h)\beta\) also commutes with \(\theta\),
the equality \(\beta\theta=\theta\beta\) gives
\(\operatorname{Ad}(h)=\operatorname{Ad}(\theta(h))\).
The center of \(E\) is trivial, so \(h=\theta(h)=1\).
This proves the stated rigidity. An inner conjugation is included
because it changes \(\beta\) by the inner coefficient
\(h\beta(h)^{-1}\). No outer-class classification follows.

\emph{5. The exact noncusp branch and its pro-odd kernel.}
The action \(T\) satisfies the full hypotheses of C.8:
it inverts the actual commutator cusp, and its homology is
\(\operatorname{diag}(-1,1)\), the identity modulo two.
There are precisely two section classes. The coefficient
\(\Phi^{-1}\in E'\) places \(\beta\) in the cusp class,
while \(X\beta\) has the other, odd \(X\)-label.
C.11 and C.12 apply to \(B=E^b\), and give
\[
 B=M\rtimes\overline{\langle a_B\rangle},\quad
 (a_B)_{\rm ab}=[Y],\quad
 \chi=\exp_Y:B\twoheadrightarrow\widehat{\mathbb Z},\quad
 M=B\cap E'\text{ pro-odd}.                                 \tag{A.130.9}
\]
At this stage uniqueness of \(a_B\) is not assumed.
Because \(b,t\) commute, \(t\) preserves \(B\).
Its abelian action is \(-I\), so it preserves \(M\) and
negates \(\chi\). Also \(M^t=1\).

\emph{6. An explicit fixed lift and elimination of the whole kernel.}
Let \(J\) be open normal and \(t\)-stable in \(M\), and suppose
\(xJ\) is fixed. Then \(n=x^{-1}t(x)\in J\) satisfies
\(t(n)=n^{-1}\). The procyclic group generated by \(n\) is
pro-odd, so its unique half-powers exist. For \(y=n^{-1/2}\),
\[
                    y^{-1}t(y)=n,\qquad
                    t(xy^{-1})=xy^{-1}.                     \tag{A.130.10}
\]
Thus every fixed coset lifts to a fixed element of \(M\).
It follows that \((M/J)^t=1\).
Such \(J\) are cofinal, by intersecting an open normal subgroup
with its \(t\)-image. On each finite quotient the injective
map \(g\mapsto g^{-1}t(g)\) is bijective; \(t\) inverts every
value of that map, hence every element. Inversion is then an
automorphism, so that finite quotient is abelian.
Consequently \(M\) is abelian and \(t(m)=m^{-1}\) for all
\(m\in M\). This argument supplies the full coprime passage;
it requires no additional theorem about arbitrary profinite
fixed-point-free involutions.

Write \(t(a_B)=m_0a_B^{-1}\), with \(m_0\in M\), and put
\(\lambda=\operatorname{Ad}(a_B)|_M\). Applying \(t\) to
\(a_Bma_B^{-1}\) gives
\[
 \lambda(m^{-1})
   =m_0a_B^{-1}m^{-1}a_Bm_0^{-1}
   =\lambda^{-1}(m^{-1}).
\]
Thus \(\lambda^2=1\), so \(a_B^2\) centralizes \(M\).
The full primitive-homology statement C.5 applies to this
nonzero integral power and gives
\(C_E(a_B^2)=\overline{\langle a_B\rangle}\).
Intersecting with \(\ker\chi\) proves \(M=1\).
This proves (A.130.4) and makes its homology-normalized
generator unique. C.8 supplies the conjugacy assertion for
the other noncusp sections.
\end{proof}

\paragraph{Canonical and remaining scope.}
For the original canonical norm word at its established
\(\widehat{\mathrm{GT}}_0\) scope,
\[
 \begin{gathered}
 U=u(X,K),\quad V=u(z,K),\\
 T(U)=U\Phi^{-1},\quad T(V)=K^{-1}VQ^{-1}K,\\
 \beta=\operatorname{Ad}(U^{-1})T\operatorname{Ad}(U).
 \end{gathered}
\]
The second identity follows from \(V=j(U)\) and (A.130.2).
No canonical-word pentagon is used. If
\(Y_0=X^{-1}a_BX\), then
\(\beta(Y_0)=XY_0X^{-1}\) and
\(\theta(Y_0)=XY_0^{-1}X^{-1}\).
The endomorphism \(X\mapsto X,\ Y\mapsto Y_0\) has identity
abelianization; its full surjectivity is not inferred.
The original boundary equality and the three-strand parameter
\(b_0\) remain separate. In particular no full unit assertion
for the canonical \(d\) is extended from full GT to GT0 here.
The complete preserved research sources are
\path{work_A/actual_elliptic_commuting_reflections37.md} and
\path{work_A/actual_GT0_two_strand_commuting_reflections37.md}.
